% vim: ai sw=2

\documentclass[xcolor=dvipsnames,handout]{beamer}

\let\Oldemph\emph
\renewcommand{\emph}[1]{\textcolor{blue}{\Oldemph{#1}}}
\newcommand{\heading}[1]{\textbf{\textcolor{BurntOrange}{#1}}}
\newcommand{\header}[1]{\textbf{\texttt{\textcolor{WildStrawberry}{#1}}}}
\newcommand{\doc}[1]{\textbf{\texttt{\textcolor{OliveGreen}{#1}}}}

% Copyright 2004 by Till Tantau <tantau@users.sourceforge.net>.
%
% In principle, this file can be redistributed and/or modified under
% the terms of the GNU Public License, version 2.
%
% However, this file is supposed to be a template to be modified
% for your own needs. For this reason, if you use this file as a
% template and not specifically distribute it as part of a another
% package/program, I grant the extra permission to freely copy and
% modify this file as you see fit and even to delete this copyright
% notice.


\mode<presentation>
{
%  \usetheme{Warsaw}
\usetheme{Madrid}

  \setbeamercovered{transparent}
}

\usepackage[czech]{babel}
\usepackage[utf8]{inputenc}
\usepackage{times}
\usepackage[IL2]{fontenc}

\usepackage{graphicx}
\usepackage[final]{listings}
\usepackage{multirow}

\lstset{tabsize=2,numbers=none,showstringspaces=false,breaklines=true,
	breakatwhitespace=true,escapeinside={(*@}{@*)},
	basicstyle=\scriptsize,keywordstyle=\bfseries\color{OliveGreen},
	commentstyle=\color{blue}}

\title[PB173/01] % (optional, use only with long paper titles)
{PB173 -- Ovladače jádra -- Linux}

\subtitle{IX.}

\author[J. Slabý]{Jiří~Slabý}

\institute[ITI] % (optional, but mostly needed)
{ITI, Fakulta Informatiky}

\date{23. 11. 2010}

\subject{}
% This is only inserted into the PDF information catalog. Can be left
% out.



% If you have a file called "university-logo-filename.xxx", where xxx
% is a graphic format that can be processed by latex or pdflatex,
% resp., then you can add a logo as follows:

% \pgfdeclareimage[height=0.5cm]{university-logo}{university-logo-filename}
% \logo{\pgfuseimage{university-logo}}



% Delete this, if you do not want the table of contents to pop up at
% the beginning of each subsection:
% \AtBeginSubsection[]
% {
%   \begin{frame}<beamer>
%     \frametitle{Outline}
%     \tableofcontents[currentsection,currentsubsection]
%   \end{frame}
% }


% If you wish to uncover everything in a step-wise fashion, uncomment
% the following command:

%\beamerdefaultoverlayspecification{<+->}

\begin{document}

\begin{frame}
  \titlepage
\end{frame}

%%%%%%%%%%%%%%%%%%%%%%%

\begin{frame}
  \frametitle{Paměť jinak I.}

  \heading{LDD3 kap. 15 (zastaralá)}

  \bigskip

  \begin{itemize}
  \item I. (dnes): Mapování paměti jádra (mmap)
  \item II. (příště): Přímý přístup do paměti (DMA)
  \end{itemize}

\end{frame}

\begin{frame}
  \frametitle{Mapování paměti jádra}

  \heading{Předání dat do/z procesu}
  \begin{itemize}
  \item Známe: \texttt{read}, \texttt{write}, \texttt{ioctl}, \ldots
  \item U všeho nutné kopírování dat
  \begin{itemize}
  \item \texttt{copy\_\{from,to\}\_user} apod.
  \end{itemize}
  \end{itemize}

  \pause
  \medskip

  \heading{Mapování paměti}
  \begin{itemize}
  \item Namapování stránek do procesu
  \item Proces používá kus stejné paměti jako jádro
  \item Syscall: \texttt{mmap}
  \end{itemize}

\end{frame}

\begin{frame}[fragile]
  \frametitle{\texttt{mmap}}

  \begin{itemize}
  \item Jeden člen v \texttt{struct file\_operations}
  \begin{itemize}
  \item \texttt{\footnotesize int mmap(struct file *filp,
    struct vm\_area\_struct *vma)}
  \end{itemize}
  \item V userspace odpovídá volání \texttt{mmap}
  \begin{itemize}
  \item \texttt{\footnotesize void *mmap(void *addr, size\_t len,
    int prot, int flags, int fd, off\_t off)}
  \end{itemize}
  \end{itemize}

  \begin{lstlisting}[language=C]
struct vm_area_struct {
  unsigned long vm_start; /* addr or random when addr is NULL */
  unsigned long vm_end;   /* vm_end = vm_start+len */
  unsigned long vm_pgoff; /* vm_pgoff = off/PAGE_SIZE */
  unsigned long vm_flags; /* vm_flags = encoded(flags|prot) */
  pgprot_t vm_page_prot;  /* only for remap_* functions */
  ...
  const struct vm_operations_struct *vm_ops; /* later ... */
  void *vm_private_data;
}
  \end{lstlisting}

  \pause

  Je třeba namapovat stránky mezi \texttt{vm\_start} a \texttt{vm\_end}. \\
  Ale také ověřit privilegia (čtení, zápis, spuštění).

\end{frame}

\begin{frame}[fragile]
  \frametitle{Základní \texttt{mmap} funkce}

  \heading{API}
  \begin{itemize}
  \item \header{linux/mm.h}, \texttt{struct vm\_area\_struct}
  \item \texttt{\_\_get\_free\_page*} $\Rightarrow$ \texttt{remap\_pfn\_range}
  \item \texttt{vmalloc\_user} $\Rightarrow$ \texttt{remap\_vmalloc\_range}
  \end{itemize}

  \begin{lstlisting}[language=C]
int my_init(void)
{
  mem = __get_free_pages(GFP_KERNEL, 2);
  /* mem = vmalloc_user(PAGE_SIZE); */
}
...
int my_mmap(struct file *filp, struct vm_area_struct *vma)
{
  if ((vma->vm_flags & (VM_WRITE | VM_READ)) != VM_READ)
    return -EINVAL;
  return remap_pfn_range(vma, vma->vm_start, virt_to_phys(mem) >> PAGE_SHIFT, 4 * PAGE_SIZE, vma->vm_page_prot);
  /* return remap_vmalloc_range(vma, mem, 0); */
}
  \end{lstlisting}

\end{frame}

\begin{frame}
  \frametitle{Úkol}

  \heading{Přemapování prostorů}
  \begin{enumerate}
  \item Alokovat souvislé 2 a 2 stránky (vmalloc a page alokátory)
  \item Zapsat na každou stránku libovolný různý řetězec
  \item Přemapovat stránky v \texttt{mmap}
  \begin{itemize}
  \item Jedna dvojice R/W, druhá RO
  \item $0 \leq \text{\texttt{pgoff}} < 2$ $\Rightarrow$ jedno mapování
  \item $2 \leq \text{\texttt{pgoff}} < 4$ $\Rightarrow$ druhé mapování
  \end{itemize}
  \item Z userspace 2$\times$ \texttt{mmap} s off 0 a 2*4096
  \begin{itemize}
  \item V pb173/09
  \end{itemize}
  \end{enumerate}

  \medskip

  \heading{API znovu}
  \begin{itemize}
  \item \header{linux/mm.h}, \texttt{struct vm\_area\_struct}
  \item \texttt{\footnotesize int mmap(struct file *filp,
    struct vm\_area\_struct *vma)}
  \item \texttt{\_\_get\_free\_page*} $\Rightarrow$ \texttt{remap\_pfn\_range}
  \item \texttt{vmalloc\_user} $\Rightarrow$ \texttt{remap\_vmalloc\_range}
  \end{itemize}

\end{frame}

\begin{frame}
  \frametitle{\texttt{mmap} po stránkách}

  \heading{Přemapování roztroušených stránek}
  \begin{itemize}
  \item Přes \texttt{remap\_pfn\_range} obtížně
  \item Při výpadcích stránek se mapují tyto stránky jednotlivě
  \begin{itemize}
  \item Každý ovladač má ,,page fault handler''
  \item Stará jádra: \texttt{nopage}
  \item Nová jádra: \texttt{fault}
  \end{itemize}
  \item V \texttt{mmap}/\texttt{nopage}/\texttt{fault} je třeba zkontrolovat
    rozsahy a velikosti
  \begin{itemize}
  \item Před tím to dělaly \texttt{remap\_*\_range} funkce
  \end{itemize}
  \item \texttt{vma->vm\_ops}
  \begin{itemize}
  \item Struktura ukazující na háčky (včetně \texttt{nopage}/\texttt{fault})
  \end{itemize}
  \item \texttt{vma->vm\_private\_data}
  \begin{itemize}
  \item Pro naše potřeby
  \item K předání informací z \texttt{mmap} do \texttt{vm\_ops}
  \end{itemize}
  \end{itemize}

\end{frame}

\begin{frame}[fragile]
  \frametitle{\texttt{mmap} po stránkách (ponovu)}

  \heading{Nové API}
  \begin{lstlisting}[language=C]
struct vm_operations_struct {
  void (*open)(struct vm_area_struct *vma);
  void (*close)(struct vm_area_struct *vma);
  int (*fault)(struct vm_area_struct *vma, struct vm_fault *vmf);
}
int my_fault(struct vm_area_struct *vma, struct vm_fault *vmf)
{ /* my_data = vma->vm_private_data; */
        unsigned long offset = vmf->pgoff << PAGE_SHIFT;
        struct page *page;
        page = my_find_page(offset);
        if (!page)
                return VM_FAULT_SIGBUS;
        get_page(page);
        vmf->page = page;
        return 0;
}
int my_mmap(struct file *filp, struct vm_area_struct *vma)
{ /* don't forget to check ranges */
  vma->vm_ops = &my_vm_ops;
  vma->vm_private_data = my_data;
  return 0;
}
  \end{lstlisting}

\end{frame}

\begin{frame}[fragile]
  \frametitle{\texttt{mmap} po stránkách (postaru)}

  \heading{Staré API}
  \begin{lstlisting}[language=C]
struct vm_operations_struct {
  ...
  struct page (*nopage)(struct vm_area_struct *vma, unsigned long address, int *type)
}
struct page *my_nopage(struct vm_area_struct *vma, unsigned long address, int *type)
{ /* my_data = vma->vm_private_data; */
  unsigned long offset = vma->pgoff << PAGE_SHIFT;
  struct page *page;

  offset += address - vma->vm_start;
  page = my_find_page(offset);
  if (!page)
    return NOPAGE_SIGBUS;
  get_page(page);

  if (type)
    *type = VM_FAULT_MINOR;

  return page;
}
  \end{lstlisting}

\end{frame}

\begin{frame}
  \frametitle{Úkol}

  \heading{Mapování roztroušených stránek} (domácí)
  \begin{enumerate}
  \item Předchozí příklad rozšiřte
  \item Místo alokací dvojstránky alokovat 2 samostatné stránky
  \begin{itemize}
  \item \texttt{\_\_get\_free\_pages(1)} $\rightarrow$ 2$\times$
    \texttt{\_\_get\_free\_page}
  \end{itemize}
  \item Přemapovat stránky v \texttt{mmap}
  \begin{itemize}
  \item \texttt{remap\_pfn\_range} $\rightarrow$ \texttt{vma->vm\_ops->nopage}
  \end{itemize}
  \item Userspace program stejný
  \end{enumerate}

  \medskip

  \heading{Potřebné podkroky}
  \begin{itemize}
  \item Definice \texttt{\footnotesize struct page *nopage(struct
    vm\_area\_struct *vma, unsigned long address, int *type)}
  \item Kontrola rozsahů v \texttt{mmap} (\texttt{end-start < 2*PAGE\_SIZE}
    apod.)
  \item Definice \texttt{vm\_ops} a přiřazení do \texttt{vma->vm\_ops}
  \end{itemize}

\end{frame}

\end{document}

